% GENERATED FILE. Edit canonical lessons, then run npm run generate:books.
% canonical-source-hash: fnv1a64:91273170e30c1e37
% canonical-lessons: TE-C273-vaddi, TE-C273-billu, TE-C273-taggimpu, TE-C273-sarukulu, TE-C273-konugolu

\chapter{Interest, Bill, Discount, Groceries, Purchase}
\label{ch:te-interest-bill-discount-groceries-purchase}
\hlchaptermodality{273}

\begin{chapteropening}
\textbf{By the end of this chapter:} \emph{I can say interest, a bill, a discount, groceries and a purchase.}

Complete the last lesson of chapter 273: I can say interest, a bill, a discount, groceries and a purchase.
\end{chapteropening}

\section[\texorpdfstring{\te{వడ్డీ} (vaḍḍī)}{వడ్డీ (vaḍḍī)}]{\te{వడ్డీ} (vaḍḍī) — interest (on money)}

\label{lesson:TE-C273-vaddi}



\begin{quote}
Before the new one: say the Telugu for a stem, then the Telugu for a wing.
\end{quote}

\subsection*{You'll want to know: \te{వడ్డీ}}
\textbf{\te{వడ్డీ}} — \emph{vaḍḍī} — \textquotedblleft{}interest (on money)\textquotedblright{}.

\begin{cousinweb}
\textbf{\te{వడ్డీ}} (\emph{vaḍḍī}) comes from Sanskrit \emph{vṛddhi}, \textquotedblleft{}increase\textquotedblright{}, through Prakrit \emph{vaḍḍhi} (CDIAL 12076): a tadbhava. CDIAL lists words for \textquotedblleft{}interest\textquotedblright{} from it in Hindi, Bengali and Old Sinhala.
\end{cousinweb}

\begin{grammarlens}[title={What you've built}]
One more word for this chapter.
\end{grammarlens}

\subsection*{Your turn}
\begin{itemize}
\raggedright
  \item \emph{Say it:} \emph{vaḍḍī}
  \item \emph{Say it:} \emph{vaḍḍī}, once more
  \item \emph{Say it:} say \emph{vaḍḍī}
  \item \emph{Recall:} say the Telugu for a stem, then the Telugu for a wing, then say \emph{vaḍḍī} again
\end{itemize}

\begin{tcolorbox}[breakable,colback=teal!4,colframe=teal!35!black,title={Before you move on}]
What does \te{వడ్డీ} mean? (\textquotedblleft{}Interest (on money)\textquotedblright{}.) Say it once more.
\end{tcolorbox}
\section[\texorpdfstring{\te{బిల్లు} (billu)}{బిల్లు (billu)}]{\te{బిల్లు} (billu) — a bill}

\label{lesson:TE-C273-billu}



\begin{quote}
Before the new one: say the Telugu for a wing, then the Telugu for interest.
\end{quote}

\subsection*{You'll want to know: \te{బిల్లు}}
\textbf{\te{బిల్లు}} — \emph{billu} — \textquotedblleft{}a bill\textquotedblright{}.

\begin{cousinweb}
\textbf{\te{బిల్లు}} (\emph{billu}) is a loan from English \emph{bill}.
\end{cousinweb}

\begin{grammarlens}[title={What you've built}]
One more word for this chapter.
\end{grammarlens}

\subsection*{Your turn}
\begin{itemize}
\raggedright
  \item \emph{Say it:} \emph{billu}
  \item \emph{Say it:} \emph{billu}, once more
  \item \emph{Say it:} say \emph{billu}
  \item \emph{Recall:} say the Telugu for a wing, then the Telugu for interest, then say \emph{billu} again
\end{itemize}

\begin{tcolorbox}[breakable,colback=teal!4,colframe=teal!35!black,title={Before you move on}]
What does \te{బిల్లు} mean? (\textquotedblleft{}A bill\textquotedblright{}.) Say it once more.
\end{tcolorbox}
\section[\texorpdfstring{\te{తగ్గింపు} (taggimpu)}{తగ్గింపు (taggimpu)}]{\te{తగ్గింపు} (taggimpu) — a discount}

\label{lesson:TE-C273-taggimpu}



\begin{quote}
Before the new one: say the Telugu for interest, then the Telugu for a bill.
\end{quote}

\subsection*{You'll want to know: \te{తగ్గింపు}}
\textbf{\te{తగ్గింపు}} — \emph{taggimpu} — \textquotedblleft{}a discount\textquotedblright{}.

\begin{cousinweb}
\textbf{\te{తగ్గింపు}} (\emph{taggimpu}) is native Dravidian, built on \emph{taggu}, \textquotedblleft{}to decrease\textquotedblright{}, which DEDR 3178 lists with Tamil \emph{tāḻ}, \textquotedblleft{}to decrease\textquotedblright{}.
\end{cousinweb}

\begin{grammarlens}[title={What you've built}]
One more word for this chapter.
\end{grammarlens}

\subsection*{Your turn}
\begin{itemize}
\raggedright
  \item \emph{Say it:} \emph{taggimpu}
  \item \emph{Say it:} \emph{taggimpu}, once more
  \item \emph{Say it:} say \emph{taggimpu}
  \item \emph{Recall:} say the Telugu for interest, then the Telugu for a bill, then say \emph{taggimpu} again
\end{itemize}

\begin{tcolorbox}[breakable,colback=teal!4,colframe=teal!35!black,title={Before you move on}]
What does \te{తగ్గింపు} mean? (\textquotedblleft{}A discount\textquotedblright{}.) Say it once more.
\end{tcolorbox}
\section[\texorpdfstring{\te{సరుకులు} (sarukulu)}{సరుకులు (sarukulu)}]{\te{సరుకులు} (sarukulu) — groceries, goods}

\label{lesson:TE-C273-sarukulu}



\begin{quote}
Before the new one: say the Telugu for a bill, then the Telugu for a discount.
\end{quote}

\subsection*{You'll want to know: \te{సరుకులు}}
\textbf{\te{సరుకులు}} — \emph{sarukulu} — \textquotedblleft{}groceries, goods\textquotedblright{}.

\begin{cousinweb}
\textbf{\te{సరుకులు}} (\emph{sarukulu}) is native Dravidian. DEDR 2353 lists \emph{saraku}, \textquotedblleft{}article, commodity\textquotedblright{}, with Tamil \emph{carakku} and Kannada \emph{saraku}, \textquotedblleft{}goods\textquotedblright{}; \textbf{\te{సరుకులు}} (\emph{sarukulu}) is its plural.
\end{cousinweb}

\begin{grammarlens}[title={What you've built}]
One more word for this chapter.
\end{grammarlens}

\subsection*{Your turn}
\begin{itemize}
\raggedright
  \item \emph{Say it:} \emph{sarukulu}
  \item \emph{Say it:} \emph{sarukulu}, once more
  \item \emph{Say it:} say \emph{sarukulu}
  \item \emph{Recall:} say the Telugu for a bill, then the Telugu for a discount, then say \emph{sarukulu} again
\end{itemize}

\begin{tcolorbox}[breakable,colback=teal!4,colframe=teal!35!black,title={Before you move on}]
What does \te{సరుకులు} mean? (\textquotedblleft{}Groceries, goods\textquotedblright{}.) Say it once more.
\end{tcolorbox}
\section[\texorpdfstring{\te{కొనుగోలు} (konugōlu)}{కొనుగోలు (konugōlu)}]{\te{కొనుగోలు} (konugōlu) — a purchase}

\label{lesson:TE-C273-konugolu}



\begin{quote}
Before the new one: say the Telugu for a discount, then the Telugu for groceries.
\end{quote}

\subsection*{You'll want to know: \te{కొనుగోలు}}
\textbf{\te{కొనుగోలు}} — \emph{konugōlu} — \textquotedblleft{}a purchase\textquotedblright{}.

\begin{cousinweb}
\textbf{\te{కొనుగోలు}} (\emph{konugōlu}) is built on native \emph{konu}, \textquotedblleft{}to buy, take\textquotedblright{} (DEDR 2151), one word with Tamil \emph{koḷ}, \textquotedblleft{}to receive, buy\textquotedblright{}.
\end{cousinweb}

\begin{grammarlens}[title={What you've built}]
One more word for this chapter.
\end{grammarlens}

\subsection*{Your turn}
\begin{itemize}
\raggedright
  \item \emph{Say it:} \emph{konugōlu}
  \item \emph{Say it:} \emph{konugōlu}, once more
  \item \emph{Say it:} say \emph{konugōlu}
  \item \emph{Recall:} say the Telugu for a discount, then the Telugu for groceries, then say \emph{konugōlu} again
\end{itemize}

\begin{tcolorbox}[breakable,colback=teal!4,colframe=teal!35!black,title={Before you move on}]
What does \te{కొనుగోలు} mean? (\textquotedblleft{}A purchase\textquotedblright{}.) Say it once more.
\end{tcolorbox}
